\documentclass[../../main.tex]{subfiles}
\begin{document}

{\bf [解]}

題意の領域$A,B$の境界を与える楕円を$P,Q$とするとこれらの方程式は
\begin{align}
  \begin{cases}
    P: \Bigl(\dfrac{x}{2}\Bigr)^2 + y^2 = 1, \\
    Q: \Bigl(\dfrac{x}{2}\Bigr)^2 + \Bigl(y - \dfrac{\sqrt{6}-\sqrt{2}}{2}\Bigr)^2 = 1
  \end{cases}
\end{align}
である．$P,Q$の交点のうち$x>0$にあるものを$T$,$x<0$にあるものを$S$とする．
題意より$P, Q$ の囲む領域の面積が $M$ である．$M$の面積$S_M$を求めよう．

\begin{figure}[htb]
  \centering
  \begin{tikzpicture}[scale=2.0, >=stealth, every node/.style={font=\small}]

    % --- 1. パラメータ定義 ---
    % 楕円半径: a = 2.0, b = 1.0
    % y方向シフト: dy = (sqrt(6) - sqrt(2)) / 2 ≈ 0.5176
    % 交点の y 座標: y0 = dy / 2 ≈ 0.2588
    % 交点の x 座標: x0 = 2 * sqrt(1 - y0^2) ≈ 2 * 0.9659 = 1.9319
    \def\aEll{2.0}
    \def\bEll{1.0}
    \pgfmathsetmacro{\dyVal}{(sqrt(6) - sqrt(2)) / 2.0}
    \pgfmathsetmacro{\yInt}{\dyVal / 2.0}
    \pgfmathsetmacro{\xInt}{2.0 * sqrt(1.0 - \yInt*\yInt)}

    \coordinate (O) at (0, 0);
    \coordinate (T) at (\xInt, \yInt);
    \coordinate (S) at (-\xInt, \yInt);

    % --- 2. 共有領域 M のハッチング ---
    \begin{scope}
      \clip (0, 0) ellipse ({\aEll} and {\bEll});
      \fill[pattern=north east lines, pattern color=blue!65, opacity=0.75]
      (0, \dyVal) ellipse ({\aEll} and {\bEll});
    \end{scope}

    % --- 3. 座標軸 ---
    \draw[->, thick] (-2.5, 0) -- (2.6, 0) node[right] {$x$};
    \draw[->, thick] (0, -1.3) -- (0, 1.9) node[above] {$y$};
    \node[below left=2pt] at (O) {$O$};

    % --- 4. 楕円の輪郭線 ---
    % 楕円 P (中心 (0,0))
    \draw[very thick, blue!85!black] (0, 0) ellipse ({\aEll} and {\bEll});
    \node[blue!90!black, font=\footnotesize, below left] at (-1.4, -0.7) {$P$};

    % 楕円 Q (中心 (0, dy))
    \draw[very thick, teal!85!black] (0, \dyVal) ellipse ({\aEll} and {\bEll});
    \node[teal!90!black, font=\footnotesize, above left] at (-1.4, {0.7 + \dyVal}) {$Q$};

    % --- 5. 交点 S, T および補助線 ---
    \draw[thick, dashed, gray!60] (-\xInt, \yInt) -- (\xInt, \yInt);
    \draw[thick, dashed, gray!60] (\xInt, 0) -- (T);
    \draw[thick, dashed, gray!60] (-\xInt, 0) -- (S);

    \fill[red!85!black] (T) circle (0.7pt) node[above right=1pt, red!85!black] {$T$};
    \fill[red!85!black] (S) circle (0.7pt) node[above left=1pt, red!85!black] {$S$};

    % y 軸上の交点 y 座標目盛り
    \draw[gray!80!black] (-0.03, \yInt) -- (0.03, \yInt);
    \node[left=2pt, font=\scriptsize] at (0, \yInt) {$\frac{\sqrt{6}-\sqrt{2}}{4}$};
    % \node[below=2pt, font=\scriptsize] at (\xInt, 0) {$\frac{\sqrt{6}+\sqrt{2}}{4}$};

    % 楕円 Q の中心
    \fill (0, \dyVal) circle (0.5pt);
    \node[right=2pt, font=\scriptsize, teal!85!black] at (0, \dyVal) {$\frac{\sqrt{6}-\sqrt{2}}{2}$};

    % --- 6. 領域 M のラベル ---
    \node[draw=blue!85!black, rounded corners=2pt, inner sep=2.5pt, font=\footnotesize\bfseries, blue!90!black]
    at (1, \yInt) {$M$};

  \end{tikzpicture}
  \caption{楕円 $P, Q$ とその交点 $S, T$，および共有領域 $M$（斜線部）}
  \label{fig:ellipses_region_M}
\end{figure}

そこで $X = \dfrac{x}{2}, Y = y$ なる変換を行い，変換された各点や図形を$'$をつけて$P, Q$ を $P', Q'$ のように表す．すると$P,Q$は変換後に円になる：
\begin{align}
  \begin{cases}
    P' : X^2 + Y^2 = 1 \\
    Q' : X^2 + \Bigl(Y - \dfrac{\sqrt{6}-\sqrt{2}}{2}\Bigr)^2 = 1
  \end{cases} \label{eq:1}
\end{align}
また，この変換で $M'$ の面積$S_{M'}$は $M$ の面積の $\dfrac{1}{2}$ 倍にあたる：
\begin{align}
  S_{M} = 2S_{M'} \label{eq:2}
\end{align}

\begin{figure}[htb]
  \centering
  \begin{tikzpicture}[scale=3, >=stealth, every node/.style={font=\small}]

    % --- 1. パラメータ定義 ---
    % 円の半径 R = 1.0
    % dy = (sqrt(6) - sqrt(2)) / 2 ≈ 0.5176
    % 交点 y 座標: Y0 = dy / 2 ≈ 0.2588 = sin(15°)
    % 交点 x 座標: X0 = cos(15°) = (sqrt(6) + sqrt(2)) / 4 ≈ 0.9659
    \def\rCirc{1.0}
    \pgfmathsetmacro{\dyVal}{(sqrt(6) - sqrt(2)) / 2.0}
    \pgfmathsetmacro{\yInt}{\dyVal / 2.0}
    \pgfmathsetmacro{\xInt}{sqrt(1.0 - \yInt*\yInt)}

    \coordinate (O)  at (0, 0);
    \coordinate (T)  at (\xInt, \yInt);
    \coordinate (S)  at (-\xInt, \yInt);
    \coordinate (OQ) at (0, \dyVal);

    % --- 2. 共有領域 M' のハッチング ---
    \begin{scope}
      \clip (O) circle (\rCirc);
      \fill[pattern=north east lines, pattern color=blue!65, opacity=0.75]
      (OQ) circle (\rCirc);
    \end{scope}

    % --- 3. 座標軸 ---
    \draw[->, thick] (-1.5, 0) -- (1.6, 0) node[right] {$X$};
    \draw[->, thick] (0, -1.3) -- (0, 1.9) node[above] {$Y$};
    \node[below left=2pt] at (O) {$O$};

    % --- 4. 2つの円 P', Q' ---
    % 円 P' (中心 O)
    \draw[very thick, blue!85!black] (O) circle (\rCirc);
    \node[blue!90!black, font=\footnotesize, below left] at (-0.75, -0.75) {$P'$};

    % 円 Q' (中心 (0, dy))
    \draw[very thick, teal!85!black] (OQ) circle (\rCirc);
    \node[teal!90!black, font=\footnotesize, above left] at (-0.75, {0.75 + \dyVal}) {$Q'$};

    % --- 5. 動径と弦 ST ---
    % 動径 OT' および OS'
    \draw[thick, black] (O) -- (T);
    \draw[thick, black] (O) -- (S);

    % 弦 S'T' (交差線 y = Y0)
    \draw[thick, dashed, gray!70] (-\xInt,0) -- (S) -- (T) -- (\xInt,0);

    % --- 6. 偏角 theta = pi/12 の表示 ---
    \draw[purple!85!black, thick] (0.35, 0) arc (0:15:0.35);
    \node[purple!90!black, font=\footnotesize] at (7.5:0.48) {$\theta$};

    % --- 7. 交点プロットとラベル ---
    \fill[red!85!black] (T) circle (0.7pt) node[above right=1pt, red!85!black] {$T'$};
    \fill[red!85!black] (S) circle (0.7pt) node[above left=1pt, red!85!black] {$S'$};

    % y 軸上の目盛り
    \draw[gray!80!black] (-0.03, \yInt) -- (0.03, \yInt);
    \node[left=2pt, font=\scriptsize] at (0, \yInt) {$\frac{\sqrt{6}-\sqrt{2}}{4}$};

    \fill (OQ) circle (0.5pt);
    \node[right=2pt, font=\scriptsize, teal!85!black] at (OQ) {$\frac{\sqrt{6}-\sqrt{2}}{2}$};
    \node[below=2pt, font=\scriptsize] at (\xInt, 0) {$\frac{\sqrt{6}+\sqrt{2}}{4}$};

    % --- 8. 領域ラベルと角度注記 ---
    \node[draw=blue!85!black, fill=white, rounded corners=2pt, inner sep=2.5pt, font=\footnotesize]
    at (0.2, \yInt) {$M'$};

    \node[draw=purple!80!black, fill=white, rounded corners=2pt, inner sep=2.5pt, font=\footnotesize, align=left]
    at (1.5, -0.6) {$\theta = \frac{\pi}{12}$ \\[2pt]
    $S_{M} = 2S_{M'}$};

  \end{tikzpicture}
  \caption{変換後の円 $P', Q'$ と交点 $S', T'$，および重なり領域 $M'$}
  \label{fig:transformed_circles_region_M_prime}
\end{figure}

さて $S', T'$ の座標は\cref{eq:1}を連立してまず$x$を消去すると
\begin{align}
  \Bigl(Y - \dfrac{\sqrt{6}-\sqrt{2}}{2}\Bigr)^2 = Y^2
  \quad \therefore
  Y = \dfrac{\sqrt{6}-\sqrt{2}}{4}
\end{align}
だから，\cref{eq:1}に代入して$X$座標を計算すると
\begin{align}
  S'\left(-\dfrac{\sqrt{6}+\sqrt{2}}{4}, \dfrac{\sqrt{6}-\sqrt{2}}{4}\right) \quad
  T'\left( \dfrac{\sqrt{6}+\sqrt{2}}{4}, \dfrac{\sqrt{6}-\sqrt{2}}{4}\right)
\end{align}
である．

図のように$X$軸と線分$O'T'$のなす角$\theta$とすると
\begin{align}
  \tan \theta
  = \dfrac{(\sqrt{6}-\sqrt{2})/4}{(\sqrt{6}+\sqrt{2})/4}
  = \dfrac{(\sqrt{6}-\sqrt{2})}{\sqrt{6}+\sqrt{2}}
  = \dfrac{(\sqrt{6}-\sqrt{2})^2}{4}
  = 2-\sqrt{3}
\end{align}
だから，$\tan$の半角公式より
\begin{align}
  \cos 2\theta
  &= \dfrac{1-\tan^2\theta}{1+\tan^2\theta} \\
  &= \dfrac{1-(7-4\sqrt{3})}{1+(7-4\sqrt{3})} \\
  &= \dfrac{-3+2\sqrt{3}}{4-2\sqrt{3}} \\
  &= \dfrac{(-3+2\sqrt{3})(4+2\sqrt{3})}{4} \\
  &= \dfrac{\sqrt{3}}{2} \\
  \therefore
  \theta &= \dfrac{\pi}{12}
\end{align}
である．

\begin{figure}[htb]
  \centering
  \begin{tikzpicture}[scale=3.8, >=stealth, every node/.style={font=\small}]

    % --- 1. パラメータ定義 ---
    % theta = 75° (5/12 pi)
    % x0 = cos(75°) = (sqrt(6) - sqrt(2)) / 4 ≈ 0.2588
    % y0 = sin(75°) = (sqrt(6) + sqrt(2)) / 4 ≈ 0.9659
    \def\angVal{75}
    \pgfmathsetmacro{\xZero}{cos(\angVal)}
    \pgfmathsetmacro{\yZero}{sin(\angVal)}

    \coordinate (O)     at (0, 0);
    \coordinate (Xlow)  at (\xZero, 0);
    \coordinate (Xhigh) at (1.0, 0);
    \coordinate (Ptop)  at (\xZero, \yZero);

    % --- 2. 積分領域 (S_M' / 4) の斜線ハッチング ---
    \fill[pattern=north east lines, pattern color=blue!65, opacity=0.75]
    (Xlow) -- (Ptop)
    arc[start angle=\angVal, end angle=0, radius=1.0]
    -- cycle;

    % --- 3. 座標軸 ---
    \draw[->, thick] (-0.2, 0) -- (1.35, 0) node[right] {$X$};
    \draw[->, thick] (0, -0.2) -- (0, 1.3) node[above] {$Y$};
    \node[below left=2pt] at (O) {$O$};

    % --- 4. 単位円弧 (第1象限) と動径 ---
    \draw[very thick, blue!85!black] (1.0, 0) arc (0:90:1.0);
    \node[blue!90!black, font=\footnotesize, above right] at (0.75, 0.75) {$X^2 + Y^2 = 1$};

    % 領域の左端 (X = x0 の境界線)
    \draw[thick, black] (Xlow) -- (Ptop);

    % 角度 5/12 pi を示す動径 OPtop
    \draw[thick, purple!85!black] (O) -- (Ptop);

    % --- 5. 角度弧と 5/12 pi のラベル ---
    \draw[thick, purple!85!black] (0.28, 0) arc (0:\angVal:0.28);
    \node[purple!90!black, font=\footnotesize] at (37.5:0.42) {$\frac{5}{12}\pi$};

    % --- 6. 補助線・目盛り ---
    % 円上の点から y 軸への投影破線
    \draw[dashed, gray!60] (0, \yZero) -- (Ptop);

    % x 軸上の積分区間
    \fill (Xlow) circle (0.6pt);
    \node[below=3pt, font=\scriptsize] at (Xlow) {$\frac{\sqrt{6}-\sqrt{2}}{4}$};
    \fill (Xhigh) circle (0.6pt) node[below=3pt, font=\scriptsize] {$1$};

    % y 軸上の目盛り
    \draw[gray!80!black] (-0.02, \yZero) -- (0.02, \yZero);
    \node[left=2pt, font=\scriptsize] at (0, \yZero) {$\frac{\sqrt{6}+\sqrt{2}}{4}$};
    \draw[gray!80!black] (-0.02, 1.0) -- (0.02, 1.0);
    \node[left=2pt, font=\scriptsize] at (0, 1.0) {$1$};

    % 境界点プロット
    \fill[red!85!black] (Ptop) circle (0.7pt);

  \end{tikzpicture}
  \caption{$\dfrac{1}{4}S_{M'}$ を与える積分領域（斜線部）}
  \label{fig:integral_region_sm_prime_quarter_angle}
\end{figure}

従って，面積$S_{M'}$は上図の斜線部の弓形面積の$4$倍の領域で
\begin{align}
  \dfrac{1}{4}S_{M'}
  &= \int_{\frac{\sqrt{6}-\sqrt{2}}{4}}^{1} \sqrt{1 - X^2} dX \\
  &= \int_{0}^{5\pi/12} \sin \theta \cdot \sin \theta d\theta \quad (X=\cos\theta) \\
  &= \int_{0}^{5\pi/12} \cos^2 \theta d\theta \\
  &= \int_{0}^{5\pi/12} \left(\dfrac{1-\cos 2\theta}{2}\right) d\theta \\
  &= \dfrac{1}{2}\left[\theta - \dfrac{1}{2}\sin 2\theta\right]_{0}^{5\pi/12} \\
  &= \dfrac{5\pi}{24} - \dfrac{1}{4}\sin \dfrac{5\pi}{6} \\
  &= \dfrac{5\pi}{24} - \dfrac{1}{8} \\
  \therefore
  S_{M'}
  &= \dfrac{5\pi}{6} - \dfrac{1}{2}
\end{align}
である．よって\cref{eq:2}に代入して
\begin{align}
  S_{M} = \dfrac{5}{3}\pi - 1
\end{align}
が求める面積である．

{\bf [別解]}

$S',T'$の座標が求まったあと，積分をせずに初等計算で弓形の面積は計算できる．
検算として活用したい．

\begin{figure}[htb]
  \centering
  \begin{tikzpicture}[scale=3.8, >=stealth, every node/.style={font=\small}]

    % --- 1. パラメータ定義 ---
    % 中心角: 75° (5/12 pi) = 90° - 15° (pi/2 - theta)
    % 点 T' の x 座標: x0 = cos(75°) = (sqrt(6) - sqrt(2)) / 4 ≈ 0.2588
    % 点 T' の y 座標: y0 = sin(75°) = (sqrt(6) + sqrt(2)) / 4 ≈ 0.9659
    \def\angVal{75}
    \pgfmathsetmacro{\xZero}{cos(\angVal)}
    \pgfmathsetmacro{\yZero}{sin(\angVal)}

    \coordinate (O)     at (0, 0);
    \coordinate (Xlow)  at (\xZero, 0);
    \coordinate (Xhigh) at (1.0, 0);
    \coordinate (Ptop)  at (\xZero, \yZero);

    % --- 2. 差し引く三角形の薄い塗りつぶし (O - Xlow - Ptop) ---
    \fill[orange!20, opacity=0.7] (O) -- (Xlow) -- (Ptop) -- cycle;

    % --- 3. 求める弓形領域 (S_M' / 4) の斜線ハッチング ---
    \fill[pattern=north east lines, pattern color=blue!65, opacity=0.75]
    (Xlow) -- (Ptop)
    arc[start angle=\angVal, end angle=0, radius=1.0]
    -- cycle;

    % --- 4. 座標軸 ---
    \draw[->, thick] (-0.2, 0) -- (1.35, 0) node[right] {$X$};
    \draw[->, thick] (0, -0.2) -- (0, 1.3) node[above] {$Y$};
    \node[below left=2pt] at (O) {$O$};

    % --- 5. 扇形の境界線 (動径と円弧) ---
    % 単位円弧
    \draw[very thick, blue!85!black] (1.0, 0) arc (0:90:1.0);
    \node[blue!90!black, font=\footnotesize, above right] at (0.75, 0.75) {$X^2 + Y^2 = 1$};

    % 動径 OPtop (扇形の上端)
    \draw[very thick, black] (O) -- (Ptop);

    % 垂線 (三角形の垂直な辺)
    \draw[thick, orange!85!black] (Xlow) -- (Ptop);

    % --- 6. 直角記号 ---
    \draw[thin, gray!70]
    ($ (Xlow) + (0, 0.06) $) --
    ($ (Xlow) + (-0.06, 0.06) $) --
    ($ (Xlow) + (-0.06, 0) $);

    % --- 7. 中心角 75° (pi/2 - theta) の表示 ---
    \draw[thick, purple!85!black] (0.28, 0) arc (0:\angVal:0.28);
    \node[purple!90!black, font=\footnotesize] at (37.5:0.43) {$\frac{\pi}{2} - \theta$};

    % --- 8. 補助線・目盛り ---
    \draw[dashed, gray!60] (0, \yZero) -- (Ptop);

    % x 軸上の目盛り
    \fill (Xlow) circle (0.6pt);
    \node[below=3pt, font=\scriptsize] at (Xlow) {$\frac{\sqrt{6}-\sqrt{2}}{4}$};
    \fill (Xhigh) circle (0.6pt) node[below=3pt, font=\scriptsize] {$1$};

    % y 軸上の目盛り
    \draw[gray!80!black] (-0.02, \yZero) -- (0.02, \yZero);
    \node[left=2pt, font=\scriptsize] at (0, \yZero) {$\frac{\sqrt{6}+\sqrt{2}}{4}$};

    % 点プロット
    \fill[red!85!black] (Ptop) circle (0.7pt);

    % 三角形内部の注釈
    \node[orange!90!black, font=\scriptsize, align=center] at (0.13, 0.38) {直角三角形};

  \end{tikzpicture}
  \caption{扇形から直角三角形を減じることによる弓形面積 $\dfrac{S_{M'}}{4}$}
  \label{fig:geometric_subtraction_sm_prime_quarter}
\end{figure}

従って図の斜線部の面積（対称性から$S_{M'}/4$）は扇形から対応する三角形を減じたもので，
\begin{align}
  \dfrac{S_{M'}}{4}
  &= \dfrac{1}{2}\left(\dfrac{\pi}{2}-\theta\right) - \dfrac{1}{2}\cdot \dfrac{\sqrt{6}+\sqrt{2}}{4} \cdot \dfrac{\sqrt{6}-\sqrt{2}}{4} \\
  &= \dfrac{5}{24}\pi - \dfrac{1}{8}
\end{align}
であり，\cref{eq:2}に代入して
\begin{align}
  S_{M} = 8 \cdot \Bigl(\dfrac{5}{24}\pi - \dfrac{1}{8}\Bigr) = \dfrac{5}{3}\pi - 1
\end{align}
である．

\end{document}
